% ============================================================
% Chapter 4: Proposed Method: Specification and Design
%   Section: Candidate Selection
% Included from chapters/04-proposed-method/proposed-method.tex
% Source: docs/decisions/selection.md; src/services/selection/.
% How the candidates are found is Section~\ref{sec:dev-scraping}.
% ============================================================

\section{Candidate Selection}
\label{sec:design-selection}

Between discovery and analysis sits a choice. Analysing an article
costs a scrape, an enrichment and one language-model call per claim ---
about three minutes per article on the production CPU --- so the
system does not analyse everything it finds. A \emph{round} lists the
candidates discovery found for up to three topic groups, ranked; a
person, a language model or both choose at most twenty; only those
become analysis jobs. Every round is recorded with what was found, what
the model proposed and why, what was sent and who chose it, which is
the data RQ1 and RQ3 are measured on. How candidates are found, and the
filters applied to links before anything is fetched, are described in
Section~\ref{sec:dev-scraping}; this section covers what happens to
them once found.

\paragraph{The mission screen.} Topic keywords say nothing about tone:
a Culture round listed election polls and celebrity interviews beside a
free music-education programme, all on topic by their words. The screen
leaves a candidate out on any of three grounds: it is not an article (a
podcast, a roundup, a quiz, a gallery, a media advisory, a promotion,
shopping), recognised by phrases in its title or summary; its headline
reports a death; or its title and summary are nearer, in the embedding
space, to one of nine descriptions of what the publication never covers
--- politics, strikes and protests, crime, accidents and disasters,
war, celebrity, sport, markets and personal finance, service pages ---
than to every one of the 23 topics, by more than a margin. The margin
was chosen by reading 527 live candidates by hand: 0.04 in general, any
amount for politics, 0.02 for the categories whose doubtful band held
almost nothing worth keeping. The three outlets that publish only
positive news are exempt, because their editors already chose. If the
inference service cannot be reached, nothing is left out and the round
says so; the screen never passes silently. What it cannot do --- tell
the housing policy of a front page from good news on the same topic ---
is left to the next steps.

\paragraph{The ranking.} A round over two or three groups finds 60 to
80 candidates, and a person reads the first twenty. Each candidate is
scored in $[0,1]$ from what the round already holds, without fetching a
page or calling a model:
\[
  s = 0.6\,(1-e)\cdot\textit{impact} \;+\; 0.6\,e\cdot\textit{topic}
      \;+\; 0.25\cdot\textit{record} \;+\; 0.15\cdot\textit{reliability},
\]
where \textit{impact} is how much nearer the candidate comes to a
description of positive impact than to the nearest off-mission one;
\textit{topic} is how well it fits the groups asked for;
\textit{record} is the share of its source's newest items the screen
kept in this round; \textit{reliability} is the source's rating; and
$e \in [0,1]$ is the \emph{topic emphasis}, the person's own trade
between positive impact and topic, 0.6 by default. The weights were
measured on 117 candidates labelled before any score existed, where the
score separated stories worth offering from the rest with an AUC of
0.78 (95\% interval 0.67 to 0.87); the default emphasis on ten labelled
one-group rounds, where it is the value at which the first ten places
hold only the group asked for (Section~\ref{sec:results-rq1}). Two
rules act on the order rather than the score: a story several outlets
carry is offered once before any story is offered twice, and a
candidate nearer another group's topics than the round's is listed
after every candidate that fits. The best twenty are shown, forty on
request.

\paragraph{The AI selection.} On request, a language model scores
every shown candidate from 0 to 10 for positive impact, from its title
and summary, against a written definition: a real change that improves
things for people or the planet beyond those in the story, backed by
something checkable; low for sport, celebrity, lifestyle, opinion,
crime, conflict and disasters. It is told the round's topics, not the
emphasis. Its picks are the best twenty at 6 or more, ticked on the
page with a reason each, and remain the editor's to change. A language
model is used here, where admission deliberately has none, because what
the publication means by positive impact is a judgement of what changed
and for whom, and the cost is one call per ten candidates rather than
one per article. A pick still goes through admission and the whole
pipeline: the selection decides what to spend analyses on, and replaces
no check.

\paragraph{The choice.} The editor keeps or changes the proposal and
sends at most twenty. The backend, not the page, records whether the
choice was the person's, the model's or both, and how many of the
model's picks were kept, dropped and added. A round can also be chosen
\emph{blind}: the editor sends a choice first and the model is run on
the round afterwards, so that its agreement with the editor is not
inflated by the editor having seen its picks (Section~\ref{sec:rq-metrics}).
